```latex
\chapter{Manuales y documentos de apoyo}

\section{Propósito del capítulo}

Este capítulo consolida los manuales, informes técnicos, anexos, videos de apoyo y documentos externos que acompañan el informe final del proyecto. Su objetivo es dejar una referencia organizada de la documentación disponible en la carpeta compartida de Google Drive denominada \textit{Documentación ANE}.

Los documentos referenciados en este capítulo no hacen parte editable del repositorio \LaTeX. Deben entenderse como material externo de consulta, soporte técnico, guía operativa, evidencia documental o material complementario para instalación, operación, administración, pruebas y mantenimiento del sistema.

Por esta razón, el presente capítulo no reproduce el contenido completo de dichos archivos, sino que identifica su nombre, propósito, relación con el proyecto, enlace de consulta y condición de documento no editable.

\section{Carpeta documental externa}

La documentación de apoyo se encuentra disponible en la siguiente carpeta de Google Drive:

\begin{itemize}
    \item \textbf{Carpeta Documentación ANE:} \url{https://drive.google.com/drive/folders/1Ne_Gym73XcDlqQFBKUjwLtbypIQ1od_t?usp=drive_link}
\end{itemize}

Los archivos almacenados en esta carpeta se consideran documentos externos no editables desde el presente informe. Cualquier actualización, reemplazo, corrección o control de versión debe realizarse directamente en la carpeta documental de Google Drive, manteniendo trazabilidad sobre los cambios.

\section{Inventario de documentos disponibles}

La Tabla~\ref{tab:documentos_apoyo_ane} presenta el inventario de documentos, manuales, informes y videos disponibles en la carpeta compartida.

\begin{longtable}{p{0.25\textwidth} p{0.35\textwidth} p{0.25\textwidth} p{0.09\textwidth}}
\caption{Manuales, informes, videos y documentos de apoyo disponibles en Google Drive}
\label{tab:documentos_apoyo_ane}\\

\hline
\textbf{Archivo} & \textbf{Propósito dentro del proyecto} & \textbf{Enlace de consulta} & \textbf{Estado} \\
\hline
\endfirsthead

\hline
\textbf{Archivo} & \textbf{Propósito dentro del proyecto} & \textbf{Enlace de consulta} & \textbf{Estado} \\
\hline
\endhead

\texttt{acceptance\_test 1.0.pdf} &
Documento de pruebas de aceptación del sistema. Debe utilizarse como referencia para validar funcionalidades, criterios de aceptación, cumplimiento de requisitos, evidencias de prueba y cierre técnico de entregables. Se relaciona principalmente con los capítulos de protocolos de prueba, validación y trazabilidad. &
\url{https://drive.google.com/drive/folders/1Ne_Gym73XcDlqQFBKUjwLtbypIQ1od_t?usp=drive_link} &
No editable \\

\hline

\texttt{Antena TDT ensamble.mp4} &
Video de apoyo para el ensamble físico de la antena TDT. Debe consultarse como material práctico para instalación, montaje, verificación física y apoyo al capítulo de hardware o cadena de recepción RF. &
\url{https://drive.google.com/drive/folders/1Ne_Gym73XcDlqQFBKUjwLtbypIQ1od_t?usp=drive_link} &
No editable \\

\hline

\texttt{Antena VHF\_UHF ensamble.mp4} &
Video de apoyo para el ensamble físico de la antena VHF/UHF. Debe utilizarse como referencia práctica para montaje, conexión, despliegue y revisión física de antenas asociadas al sistema de monitoreo. &
\url{https://drive.google.com/drive/folders/1Ne_Gym73XcDlqQFBKUjwLtbypIQ1od_t?usp=drive_link} &
No editable \\

\hline

\texttt{Informe Técnico\_ Análisis De Antenas.docx} &
Informe técnico de apoyo para el análisis de antenas, criterios de selección, comportamiento esperado, condiciones de uso, bandas de operación y consideraciones asociadas al sistema de monitoreo. Debe consultarse como soporte del capítulo de hardware, adquisición RF y diseño del sistema. &
\url{https://drive.google.com/drive/folders/1Ne_Gym73XcDlqQFBKUjwLtbypIQ1od_t?usp=drive_link} &
No editable \\

\hline

\texttt{Manual\_del\_Usuario.pdf} &
Manual orientado al usuario final u operador del sistema. Debe servir como guía para el uso cotidiano de la plataforma, consulta de información, interacción con la interfaz, revisión de resultados, visualización de mediciones y operación básica del sistema. &
\url{https://drive.google.com/drive/folders/1Ne_Gym73XcDlqQFBKUjwLtbypIQ1od_t?usp=drive_link} &
No editable \\

\hline

\texttt{Anexo.pdf} &
Documento complementario que agrupa información adicional, detalles de soporte, evidencias, material ampliado o contenido auxiliar que no se desarrolla completamente en el cuerpo principal del informe. Debe consultarse como anexo externo. &
\url{https://drive.google.com/drive/folders/1Ne_Gym73XcDlqQFBKUjwLtbypIQ1od_t?usp=drive_link} &
No editable \\

\hline

\texttt{Manual\_del\_administrador\_de\_la\_plataforma.pdf} &
Manual dirigido al administrador de la plataforma. Debe utilizarse como referencia para gestión de usuarios, configuración, administración de servicios, revisión de parámetros, permisos, mantenimiento lógico, operación administrativa y supervisión del sistema. &
\url{https://drive.google.com/drive/folders/1Ne_Gym73XcDlqQFBKUjwLtbypIQ1od_t?usp=drive_link} &
No editable \\

\hline

\texttt{Documentación\_\_soporte\_del\_sistema\_de\_monitoreo\_multiproposito.pdf} &
Documento de soporte del sistema de monitoreo multipropósito. Debe utilizarse como referencia general para entender la arquitectura, componentes, flujo de información, funcionalidades principales, criterios técnicos, elementos de soporte, operación general y alcance funcional del sistema. &
\url{https://drive.google.com/drive/folders/1Ne_Gym73XcDlqQFBKUjwLtbypIQ1od_t?usp=drive_link} &
No editable \\

\hline

\texttt{manual\_de\_instalacion.pdf} &
Guía de instalación del sistema. Debe servir como referencia para despliegue inicial, preparación del entorno, instalación de componentes, configuración básica, verificación posterior a la instalación y puesta en marcha. Se relaciona con los capítulos de hardware, software de borde y plataforma cloud. &
\url{https://drive.google.com/drive/folders/1Ne_Gym73XcDlqQFBKUjwLtbypIQ1od_t?usp=drive_link} &
No editable \\

\hline

\texttt{informe\_IQ\_Balance.pdf} &
Informe técnico relacionado con balance I/Q, calibración o análisis de señales complejas en sistemas SDR. Debe consultarse como soporte para los capítulos de adquisición SDR, procesamiento digital de señales, calidad de datos, calibración y validación técnica de mediciones. &
\url{https://drive.google.com/drive/folders/1Ne_Gym73XcDlqQFBKUjwLtbypIQ1od_t?usp=drive_link} &
No editable \\

\hline

\end{longtable}

\section{Relación de los documentos con los capítulos del informe}

Los documentos externos deben utilizarse como soporte de los capítulos técnicos del informe final. La relación sugerida es la siguiente:

\begin{itemize}
    \item \textbf{Capítulo de diseño del sistema:} debe apoyarse en el documento \texttt{Documentación\_\_soporte\_del\_sistema\_de\_monitoreo\_multiproposito.pdf}, el archivo \texttt{Anexo.pdf} y los documentos técnicos relacionados.
    
    \item \textbf{Capítulo de hardware y cadena RF:} debe apoyarse en el \texttt{Informe Técnico\_ Análisis De Antenas.docx}, en los videos \texttt{Antena TDT ensamble.mp4} y \texttt{Antena VHF\_UHF ensamble.mp4}, así como en el \texttt{manual\_de\_instalacion.pdf}.
    
    \item \textbf{Capítulo de adquisición SDR:} debe apoyarse en el \texttt{informe\_IQ\_Balance.pdf}, en la documentación de soporte del sistema de monitoreo y en los criterios técnicos relacionados con adquisición, calibración y calidad de señal.
    
    \item \textbf{Capítulo de plataforma web o cloud:} debe apoyarse en el \texttt{Manual\_del\_administrador\_de\_la\_plataforma.pdf}, el \texttt{Manual\_del\_Usuario.pdf} y la documentación de soporte del sistema de monitoreo multipropósito.
    
    \item \textbf{Capítulo de operación del sistema:} debe apoyarse principalmente en el \texttt{Manual\_del\_Usuario.pdf}.
    
    \item \textbf{Capítulo de instalación y despliegue:} debe apoyarse en el \texttt{manual\_de\_instalacion.pdf} y en los videos de ensamble de antenas.
    
    \item \textbf{Capítulo de mantenimiento y administración:} debe apoyarse en el \texttt{Manual\_del\_administrador\_de\_la\_plataforma.pdf}, el \texttt{manual\_de\_instalacion.pdf} y la documentación de soporte del sistema.
    
    \item \textbf{Capítulo de protocolos de prueba:} debe apoyarse en el archivo \texttt{acceptance\_test 1.0.pdf}, el cual debe servir como referencia para criterios de aceptación, validación funcional y evidencia de cumplimiento.
    
    \item \textbf{Anexos y material complementario:} deben apoyarse en el archivo \texttt{Anexo.pdf}.
\end{itemize}

\section{Clasificación funcional de la documentación}

Para facilitar la consulta, los documentos externos se pueden clasificar funcionalmente de la siguiente manera:

\begin{itemize}
    \item \textbf{Documentos de operación:}
    \begin{itemize}
        \item \texttt{Manual\_del\_Usuario.pdf}
    \end{itemize}
    
    \item \textbf{Documentos de administración:}
    \begin{itemize}
        \item \texttt{Manual\_del\_administrador\_de\_la\_plataforma.pdf}
    \end{itemize}
    
    \item \textbf{Documentos de instalación y despliegue:}
    \begin{itemize}
        \item \texttt{manual\_de\_instalacion.pdf}
        \item \texttt{Antena TDT ensamble.mp4}
        \item \texttt{Antena VHF\_UHF ensamble.mp4}
    \end{itemize}
    
    \item \textbf{Documentos técnicos de soporte:}
    \begin{itemize}
        \item \texttt{Informe Técnico\_ Análisis De Antenas.docx}
        \item \texttt{informe\_IQ\_Balance.pdf}
        \item \texttt{Documentación\_\_soporte\_del\_sistema\_de\_monitoreo\_multiproposito.pdf}
    \end{itemize}
    
    \item \textbf{Documentos de pruebas y aceptación:}
    \begin{itemize}
        \item \texttt{acceptance\_test 1.0.pdf}
    \end{itemize}
    
    \item \textbf{Material complementario:}
    \begin{itemize}
        \item \texttt{Anexo.pdf}
    \end{itemize}
\end{itemize}

\section{Requisitos previos de consulta e instalación}

Antes de aplicar los procedimientos descritos en los manuales externos, se recomienda verificar los siguientes aspectos:

\begin{itemize}
    \item Contar con acceso de lectura a la carpeta de Google Drive \textit{Documentación ANE}.
    \item Confirmar que los documentos consultados corresponden a la versión vigente o aceptada por el equipo del proyecto.
    \item Identificar el subsistema al que corresponde cada documento: hardware, SDR, plataforma, administración, instalación, operación, pruebas o soporte.
    \item Disponer del entorno técnico requerido para aplicar los procedimientos: nodo de borde, plataforma cloud, equipo SDR, red de comunicaciones, credenciales administrativas o interfaz de usuario.
    \item Validar que los procedimientos de instalación, administración o mantenimiento sean coherentes con el alcance real del prototipo o sistema implementado.
    \item Evitar copiar en el informe credenciales, direcciones internas, configuraciones sensibles o información operativa que deba mantenerse bajo control de acceso.
    \item Verificar que los videos de ensamble sean usados como apoyo práctico y no como sustituto de los procedimientos técnicos escritos.
\end{itemize}

\section{Instrucciones de operación cotidiana}

Para consultar y utilizar la documentación de apoyo durante la operación del sistema, se recomienda el siguiente procedimiento:

\begin{enumerate}
    \item Abrir la carpeta compartida \textit{Documentación ANE} desde el enlace indicado en este capítulo.
    
    \item Identificar el documento correspondiente al tipo de actividad que se desea realizar.
    
    \item Para operación funcional del sistema, consulta de resultados, navegación de la plataforma y uso por parte del operador, consultar el archivo:
    \begin{itemize}
        \item \texttt{Manual\_del\_Usuario.pdf}
    \end{itemize}
    
    \item Para tareas de administración, gestión de usuarios, configuración, revisión de servicios, permisos o mantenimiento lógico, consultar el archivo:
    \begin{itemize}
        \item \texttt{Manual\_del\_administrador\_de\_la\_plataforma.pdf}
    \end{itemize}
    
    \item Para instalación, despliegue inicial, preparación del entorno y puesta en marcha, consultar el archivo:
    \begin{itemize}
        \item \texttt{manual\_de\_instalacion.pdf}
    \end{itemize}
    
    \item Para montaje físico de antenas, consultar los videos:
    \begin{itemize}
        \item \texttt{Antena TDT ensamble.mp4}
        \item \texttt{Antena VHF\_UHF ensamble.mp4}
    \end{itemize}
    
    \item Para aspectos técnicos relacionados con antenas, criterios de selección, análisis RF o comportamiento esperado, consultar el archivo:
    \begin{itemize}
        \item \texttt{Informe Técnico\_ Análisis De Antenas.docx}
    \end{itemize}
    
    \item Para aspectos relacionados con adquisición SDR, señales I/Q, balance I/Q o calibración, consultar el archivo:
    \begin{itemize}
        \item \texttt{informe\_IQ\_Balance.pdf}
    \end{itemize}
    
    \item Para información general de soporte del sistema de monitoreo multipropósito, consultar el archivo:
    \begin{itemize}
        \item \texttt{Documentación\_\_soporte\_del\_sistema\_de\_monitoreo\_multiproposito.pdf}
    \end{itemize}
    
    \item Para pruebas de aceptación, criterios de validación y cierre funcional, consultar el archivo:
    \begin{itemize}
        \item \texttt{acceptance\_test 1.0.pdf}
    \end{itemize}
    
    \item Para información complementaria o material ampliado, consultar el archivo:
    \begin{itemize}
        \item \texttt{Anexo.pdf}
    \end{itemize}
    
    \item Registrar cualquier inconsistencia entre el documento externo y el informe final para revisión posterior.
\end{enumerate}

\section{Solución de problemas frecuentes}

La Tabla~\ref{tab:problemas_documentos_apoyo} resume problemas comunes asociados a la consulta o uso de los documentos externos.

\begin{longtable}{p{0.30\textwidth} p{0.32\textwidth} p{0.30\textwidth}}
\caption{Problemas frecuentes en la consulta de documentos externos}
\label{tab:problemas_documentos_apoyo}\\

\hline
\textbf{Problema} & \textbf{Causa probable} & \textbf{Acción recomendada} \\
\hline
\endfirsthead

\hline
\textbf{Problema} & \textbf{Causa probable} & \textbf{Acción recomendada} \\
\hline
\endhead

No se puede abrir la carpeta de Drive. &
El usuario no tiene permisos de lectura o el enlace fue modificado. &
Solicitar acceso al responsable de la carpeta documental o verificar que el enlace compartido siga vigente. \\

\hline

No se puede abrir un documento específico. &
El archivo puede haber sido movido, eliminado, renombrado o tener permisos restringidos. &
Ingresar primero a la carpeta principal y ubicar el documento por nombre. Si no aparece, solicitar confirmación al responsable documental. \\

\hline

No se puede reproducir un video de ensamble. &
El archivo de video puede requerir permisos adicionales, conexión estable o reproducción desde el visor de Google Drive. &
Abrir el video directamente desde la carpeta de Drive, verificar permisos de lectura y descargarlo solo si el responsable documental lo permite. \\

\hline

El nombre del documento en el informe no coincide exactamente con el archivo en Drive. &
Puede existir cambio de nombre, duplicidad, error de escritura o diferencia entre versiones. &
Verificar el nombre completo directamente en Drive y actualizar la Tabla~\ref{tab:documentos_apoyo_ane}. \\

\hline

Existen varias versiones de un mismo documento. &
Pueden haberse generado copias de trabajo, versiones preliminares o documentos duplicados. &
Usar únicamente la versión validada por el equipo del proyecto y registrar la fecha de consulta. \\

\hline

El contenido del manual externo no coincide con el sistema implementado. &
El manual puede corresponder a una versión anterior del sistema o a una configuración distinta. &
Registrar la diferencia y validar con el responsable técnico antes de aplicar el procedimiento. \\

\hline

El documento contiene configuraciones sensibles. &
Algunos manuales pueden incluir rutas, parámetros, direcciones, usuarios de ejemplo o configuraciones internas. &
No copiar información sensible en el informe final. Mantener los documentos externos con permisos controlados. \\

\hline

El enlace apunta a la carpeta general y no al archivo específico. &
No se dispone todavía del enlace individual del archivo o se decidió conservar una referencia general a la carpeta documental. &
Copiar desde Google Drive el enlace individual del archivo y reemplazarlo en la tabla de inventario documental. \\

\hline

\end{longtable}

\section{Recomendaciones para mantenimiento documental}

Para conservar la coherencia entre el informe final y los documentos externos, se recomienda:

\begin{itemize}
    \item Mantener la carpeta \textit{Documentación ANE} como repositorio documental externo de consulta.
    \item No duplicar el contenido completo de los manuales dentro del archivo \LaTeX.
    \item Conservar los documentos externos como archivos no editables desde el informe.
    \item Agregar enlaces individuales a cada documento cuando estén disponibles.
    \item Actualizar este capítulo si se agregan, eliminan, renombran o reemplazan documentos en la carpeta de Drive.
    \item Verificar periódicamente que los enlaces de consulta funcionen correctamente.
    \item Evitar reemplazos silenciosos de documentos. Si se actualiza un archivo, se recomienda conservar control de versión o fecha.
    \item Relacionar cada documento con el capítulo técnico correspondiente para mejorar la trazabilidad.
    \item No incluir en el informe credenciales, parámetros internos sensibles o información de acceso restringido.
    \item Distinguir entre documentos técnicos, manuales de usuario, manuales de administrador, videos de ensamble, pruebas de aceptación y anexos complementarios.
\end{itemize}

\section{Control de trazabilidad documental}

Para fortalecer la trazabilidad del informe, cada documento externo debe asociarse con al menos uno de los siguientes elementos:

\begin{itemize}
    \item Requisito del sistema.
    \item Subsistema técnico.
    \item Procedimiento de instalación.
    \item Procedimiento de operación.
    \item Procedimiento de administración.
    \item Procedimiento de mantenimiento.
    \item Procedimiento de prueba o aceptación.
    \item Evidencia de validación.
    \item Decisión de diseño.
    \item Material complementario o anexo.
\end{itemize}

La trazabilidad permite identificar qué documento respalda una afirmación, configuración, prueba o procedimiento incluido en el informe. Por tanto, cuando un documento externo sea utilizado como respaldo de una afirmación técnica, debe indicarse en el capítulo correspondiente mediante una referencia textual, una nota, una tabla de trazabilidad o una mención explícita al archivo consultado.

\section{Nota sobre enlaces individuales}

En la versión actual de este capítulo se referencia la carpeta principal de Google Drive para todos los documentos. Esto permite acceder al conjunto completo de archivos disponibles en la carpeta \textit{Documentación ANE}. Sin embargo, para mejorar la trazabilidad documental, se recomienda reemplazar progresivamente el enlace general por el enlace individual de cada archivo.

El procedimiento sugerido es:

\begin{enumerate}
    \item Abrir la carpeta \textit{Documentación ANE}.
    \item Hacer clic derecho sobre cada archivo.
    \item Seleccionar la opción de compartir o copiar enlace.
    \item Verificar que el permiso permita consulta de lectura.
    \item Reemplazar en la Tabla~\ref{tab:documentos_apoyo_ane} el enlace general de la carpeta por el enlace individual del archivo.
    \item Conservar el estado \textit{No editable} para indicar que el documento es una fuente externa de consulta.
\end{enumerate}

\section{Observación final}

Los manuales, informes, videos y documentos de soporte incluidos en este capítulo deben considerarse parte del ecosistema documental del proyecto, pero no sustituyen el contenido técnico principal del informe. Su función es complementar, respaldar y facilitar la consulta de información específica relacionada con instalación, operación, administración, pruebas, análisis técnico y mantenimiento del sistema de monitoreo.
```
